\section{Evidence, counterexamples, and remaining coverage}\label{sec:experiments}

\subsection{Retained baseline evidence}
Manuscript v1.0 audited the public \texttt{0.1.0a1} commit on Python 3.12.14 and x86-64 Linux. Its retained evidence includes 67 passing tests, 22,140 concrete comparisons of the original AND/OR/XOR programs at widths one through four, and an independent matrix-versus-ground comparison on 804 instances and 35,856 interface pairs. All 29 rewrite identifiers had witnesses, with 23,751 finite before/after comparisons and checks of the decreasing measure.

The finite classification checked all 512 Boolean row functions and all $3^9$ valid output tables per target, confirming sound-class counts $64,64,16$ and monotone counts $26,26,16$. Of 680 source mutations, 270 were certified, 180 refuted at width one, and 230 required fallback. These remain evidence about the pinned coordinatewise profile. Versioned logs retain the original timings; this revision does not treat them as new performance measurements.

\subsection{Broader research has several denominators}
\begin{table}[htbp]\small
\begin{tabularx}{\textwidth}{@{}l l X@{}}
\toprule
Dataset or profile & Result & Interpretation\\
\midrule
Known LLVM transcriptions & $11/11$ & A development/regression set, distinct from the original 39-program corpus.\\
GCC 15.2 specializations & $9/9$ & Frozen external transfer study over the selected nine dispatcher specializations.\\
Original Nice to Meet You solutions & $11/39$ & Complete source-bound proofs in the declared research contract.\\
Their individual components & $104/411$ & Component proofs; not additional complete programs.\\
Five Graal procedures, control & $1/5$ & Complete mask result for OR.\\
Five Graal procedures, joint & $2/5$ & Complete mask results for AND and OR.\\
\bottomrule
\end{tabularx}
\caption{These populations and proof obligations are not interchangeable.}\label{tab:coverage}
\end{table}

On the original 39 solutions, the result-observation stage increased the preceding control from $9/39$ complete programs and $87/411$ components to $11/39$ and $104/411$. Add reached $13/13$ components and UsubSat $16/16$. All 411 source components were prepared; source-to-expression preparation alone is not target verification. The other 28 programs include 23 targets without the required research target adapter and five flag-related strengthened obligations. Refutations of those unconditional strengthened obligations do not establish errors in the upstream guarded contracts.

The historical GCC study froze its backend and selection before outcomes. It certified nine SSA specializations from GCC 15.2 and replayed them in separate processes. Composition reduced final certificate states from 2344 to 1070 at unchanged coverage. Accounting for all search attempts gives 2344 versus 1474 states; the smaller certificate count is not the full search cost. Z3 also solved all 72 selected finite-width obligations. These are transfer and structural results on that selected corpus.

\subsection{Context and reconstruction evidence}
The whole-word reconstruction study retained all claims of a normalization query while replacing a 4401-node subproof by 26 recorded steps, including the universal projection lemma, applicability conditions, ground steps, and substitution. The complete proof decreased from 4905 to 530 steps. On its 18 constructed evaluation sources, the new mechanism proved 11 intended positive cases, left two positive cases unsupported, and returned unresolved on five faulty cases for which separate finite counterexamples were retained. It was a constructed evaluation, not a blind external benchmark.

The evidence also preserves cases where observation machinery increased cost without reducing a proof. The three-bit equal-kernel/different-label separator in Section~\ref{sec:observations} and the two-bit shared-prefix separator in Section~\ref{sec:wordbounds} identify concrete information that cannot be discarded. Neither is an upstream program bug.

\subsection{Upper and lower contract checks}
\begin{center}\small
\begin{tabularx}{\textwidth}{@{}lX@{}}
\toprule
Check & Recorded result\\
\midrule
Upper, all masks/bounds/zero flags at widths 1--6 & 111,972 cases agree with exact masked maximum and emptiness.\\
Upper, extracted Java helper & 114,040 calls agree with the declared model.\\
Upper, additional lower bounds at widths 1--4 & 24,174 joint carriers preserved exactly.\\
Lower, all small typed cases at widths 1--6 & 111,972 cases: 76,104 admitted by the sign contract; 35,868 explicitly unsupported.\\
Lower, extracted Java helper & 114,044 calls agree with the declared model.\\
Lean pilot & Clean build, fresh extraction build, axiom audit, three examples, and four rejected isolated corruptions.\\
\bottomrule
\end{tabularx}
\end{center}

The universal upper certificate additionally has wide mathematical specializations, including 128-, 256-, and 4096-bit cases. These are not Java runs at those widths. Upper and lower tables each describe their own experiment; their related input populations are not added together to inflate coverage. Among upper returns numerically equal to the minimum word, the record separates 37,049 empty cases from 963 genuine maxima.

The lower study retains conservative and out-of-contract behavior, including the examples following Theorem~\ref{thm:lower}. A strict-comparison source variant was rejected by the reused floor lemma even though its 2232 checked native cases did not exhibit a bug. Such rejection marks the boundary of a proof interface; it is not by itself an unsoundness verdict.

\subsection{What the SMT comparisons establish}
One historical comparison used the same stored source-derived formulas for 39 programs at widths 32, 64, and 128, with three solvers and two repeats: 117 obligations and 702 executions. Its unconditional treatment of flag-related programs was later separated from artifact and LLVM-compatible assumptions. A subsequent 32-bit study varied both that contract and the encoding of counts, with a one-second solver budget.

The latter change was substantial. Under the artifact contract, balanced encoding gave 19 UNSAT results for Z3 and 18 for Bitwuzla out of 39, with no SAT results; their prefix encodings gave 13 each. Under the historical contract, balanced encoding gave 16 and 20 SAT results respectively. These runs show why source semantics and encoding must be held explicit before a ranking is meaningful. Their historical versions, budgets, and records remain attached to the reports.

QKF supplies parametric proofs for supported interfaces; these SMT runs ask separate fixed-width questions and did not retain solver proof objects. A useful local result was an all-width QKF proof for Smax when the six recorded 128-bit SMT attempts reached their limits, while at 32 bits conventional solvers could be faster. The evidence supports further controlled comparison and identifies specific successful interfaces. It does not establish a general performance ordering or coverage of arbitrary SMT problems.

\subsection{Later interface and algebraic evaluations}
The completed full-source frontend-applicability study obtained $0/12$ accepted external methods across two Graal-specific frontends (24 admission attempts). All attempts stopped at source admission; no source-model search or target proof was reached, and no program was refuted. This deliberately selected, previously inspected sample measures applicability beyond the old frontend contract. Its final acquisition and CI record is \href{https://github.com/Len989/QKF-Certifier/pull/14}{PR14}; \artifact{research/external/frozen_v1/REPORT_RU.md} preserves the earlier provisional run and must not be read as the final status. Later extension work does not retroactively change this result. The later frozen Java external studies retain their own denominators: $1/15$ in frozen-v2 and $2/33$ in Run27. The two Run27 successes are two overloads of one body/algorithm group; the other 31 cases are capability or contract gaps, not source refutations. Later reuse of development methods is not a new holdout evaluation. The locked corpus and acquisition rules are preserved in \artifact{research/external/frozen_v2/PROTOCOL.md} and \artifact{research/external/frozen_v3/run27/README_RU.md}. These results neither replace nor increase the earlier $11/39$ program and $104/411$ component figures.

The source-interface acceptance study used 23 main development cases and five explicit controls, comparing the frozen route, the new row route, and a direct-cell control. The old route completed all 23 main cases with 15 certified and eight refuted; both new routes completed 21 with 14 certified and seven refuted, leaving two source-cap failures. Truth outcomes agreed on their shared completed population, and the new routes had identical source interfaces and target obligations. A smaller factor did not imply a smaller proof: one delayed case fell from 61 to 33 classes while its complete packet grew from 22,240 to 92,017 bytes. See \artifact{research/acceptance/REPORT_RU.md}. The later reduction and coverage mechanisms address the reported cap failures without retrospectively changing this experiment.

The pure-row evaluation comprises 16 named, 5120 exhaustive two-element, and 128 seeded mixed presentations: 5264 finite algebraic inputs, not Java programs. The ground-query conformance set contains 203 registered requests, including 195 Catalan profiles through height five and eight named variants. It compares each activation horizon against the retained independent full-space ground procedure; a separate finite rewrite search exhibits a congruence/rewrite-depth gap of $2/3$. These are correctness and representation studies described in the respective module READMEs, not evidence of a general speed advantage.

\subsection{Selecting facts and comparing equal deliveries}
The causal study in \artifact{research/causal_comparison/REPORT_RU.md} includes 32 cases, 122 scheduled case/route pairs, six preregistered unsupported contracts, and 116 executed pairs. It records all 1044 attempts. Adaptive acquisition reduces native questions on the mask and low-bit families: respectively $2/20$, $4/42$, $13/23$, and $8/36$ adaptive/eager attempts on the four named positive source cases. The complete adaptive SDK nevertheless costs more than the simpler direct route on these singletons. The experiment therefore establishes selective acquisition without establishing a complete-path speedup from it.

Reuse lowers the SDK's median build/check cost on the four $N=16$ mask/parity and all/fixed-width series in that study. These targets are distinct expressions equivalent under the same guard, not 16 semantically independent properties. Smaller final dependencies and reduced export work contribute to the result; it is not an isolated advantage of the query backend. In the local physical-phase pilot, forcing obtains a genuine value in $D\setminus S$, but direct controls already need only three physical source-cell evaluations and cost less.

The next comparison holds SDK functionality equal. It distinguishes one applicable summary, one checked module serving equivalent consumers through \texttt{assess}, and independently checkable packets for every target. QKF's ordinary-to-SDK route includes adaptation and checking costs; raw proof-only timings are reported separately. On each of Python 3.10 and 3.12, that route is cheaper than the default in all four registered singleton cases. The corresponding once-ordinary service and each-ordinary independent-packet routes are cheaper in all 13 rows of their respective contracts. These internal baselines are specified in \artifact{research/sdk_comparison/REPORT_RU.md} and \artifact{research/sdk_comparison/PR52_BASELINES.json}.

The service can build one sufficient summary and assess later equivalent consumers. This saving does not imply transfer to a different source, guard, width, or semantic contract. The independent-packet clients in that experiment did not retain a cross-call semantic or exact cache; exact-repeat conclusions therefore do not exclude a cheaper client that simply retains the checked certificate. These delivery distinctions are necessary for interpreting reuse.

\subsection{Prepared contexts and direct proof emission}
The prepared-context experiment records, on each Python version, five cases, 38 case/route pairs, and 646 attempts. Corresponding old/new deliveries and native attempts agree. Preparation reduces repeated schema parsing and reconstruction; in the mask/all-width/four-goal audit, ground parses fall from 144 to 74 and research Python calls from 1,583,455 to 1,273,278. The Python 3.12 lifecycle medians improve by 9.03--17.31\% across the five registered deliveries, with build traced-memory peaks increasing by approximately 1.6--5.5\%. These are implementation tradeoffs under unchanged proof obligations; see \artifact{research/prepared_context/REPORT_RU.md}.

Direct emission removes the intermediate positive native packet. Its experiment records six cases, 48 case/route pairs, and 816 attempts on each Python version. The 42 matched old/new comparisons, including two scoped fallbacks, preserve delivery/result bytes, native attempts, foundations, and checkpoint identities. Across ten positive environment comparisons, query/no-lemmas lifecycle medians decrease by 27.96--38.80\% relative to the prepared-context route. These are ranges of separate comparisons, not a pooled estimate or confidence interval.

\begin{table}[htbp]\small
\begin{tabularx}{\textwidth}{@{}Xrrr@{}}
\toprule
Delivery & Prepared & Direct emission & QKF ordinary\\
\midrule
Mask, one summary & 835.98 & 578.97 & 347.75\\
Mask/all widths, four-consumer service & 3482.99 & 2165.81 & 434.58\\
Parity/fixed8, four-consumer service & 3477.76 & 2156.00 & 409.86\\
Sixteen exact repeats & 2942.70 & 2119.99 & 505.14\\
Mask, four independent packets & 6478.41 & 4080.38 & 1439.30\\
Open power predicate, scoped fallback & 629.47 & 634.19 & 361.84\\
\bottomrule
\end{tabularx}
\caption{PR49 experiment: lifecycle median milliseconds on Python 3.12 within each registered job. Prepared and direct-emission columns use query/no-lemmas. The final column is QKF's simpler registered ordinary-SDK control for the same delivery contract.}\label{tab:lifecycle}
\end{table}

Table~\ref{tab:lifecycle} is drawn from \artifact{research/direct_emission/REPORT_RU.md}; raw samples and identities are retained in its \texttt{RESULTS.json}. The simpler ordinary-SDK control remains cheaper than each of the four new ordinary/query by no-lemmas/reuse configurations in all ten positive environments. Direct emission is therefore a successful optimization within QKF, while an additional advantage of the advanced query path over QKF's simpler route remains unestablished. No default switch follows from these timings.

The later experiments use seven fresh timing samples, with audit and memory runs performed separately. Lifecycle includes build, delivery encoding, independent cold checking, applicable consumer assessment, export, and explanation. Repeated application is measured separately. Traced Python build allocations are not resident memory or receiver memory; Python call counts are not machine instructions. Different runners and interpreter versions are not pooled. These are instrumented development measurements, not compiler hot-path or external-generalization results.

\subsection{Validation provenance}
The committed reports retain engine identities, registrations, raw outcomes, manifests, and replay recipes. Formal validation records are described in Section~\ref{sec:lean}. Preparing this manuscript does not itself rerun the historical searches, native executions, benchmarks, or Lean builds. Any additional release smoke checks are listed separately in the accompanying release-validation record, so a packaging check cannot be mistaken for a new research measurement.
